\chapter{Supplement B: Base Transport, Infinite Networks, and the Lattice Cluster}
\label{comp:sb}

\section{The question and the answer}
\label{comp:sb-question}

The main manuscript works at one base point \(c_*\): every frame is a point of one principal fiber,
the independent links on the edges of the agent network are the only transports, and the design is
\(\{c_*\}\). Supplement B collects what the theory would need if contexts varied over a base, if the
network became infinite, or if agents sat on a lattice. For each extension it asks which single-point
statements survive and exactly which additional data the extension consumes: a connection, a cover, a
design, a base measure, a scale identification, or a lattice realization (\SB{ch:base-introduction}).
Nothing in the main manuscript depends on it. This chapter reads it at overview depth with the scope
gate in force: it reports what the supplement proves and where it stops, and proposes nothing.

Part B-I answers the base question with kinematics. Relative frames, chosen connections, covariant
defects of cross maps, four independent notions of flatness, finite designs, a Fisher geometry pulled
back to the base, gauge-invariant durations for field histories, and a descent criterion are all
typed and proved, with counterexamples wherever an expected implication fails. No connection, base measure, or probability
law on fields is derived (\SB{open:prob-continuum-theory}), and graph links become base transports
only by declaration (\SB{hyp:geo-graph-base-transport}). Part B-II answers what Wilsonian vocabulary
requires: comparison maps that the coarse channels do not supply. A finite network at \(c_*\) carries
no asymptotics (\SB{prop:rg-finite-depth-no-asymptotics}), beta functions are scheme dependent, and on
Gaussian hierarchies attraction is proved only under a fixed primitive identification whose existence
is open (\SB{conj:grg-fixed-b-attraction}). No infinite agent network, exponent, or universality class
is constructed. Part B-III constructs one infinite-volume family completely: the square-lattice
\(\mathbb Z_2\) gauge--matter law at small plaquette coupling, with a coarse specification at every
exterior for every block size on a block-independent domain
(\SB{thm:rg-spatial-constrained-specification}), a differentiable renormalization map, and proved
volume limits of pressure, label selection, specific VFE, and Fisher densities
(\SB{cor:rg-thermo-native-label-limits}). Whether its sites are agents at \(c_*\) or base points is
open (\SB{open:rg-lattice-typing}).

The supplement does what its introduction says. Its title words promise more: no base transport of
beliefs is derived from inference, no infinite agent network is built, and the one constructed
infinite-volume family is untyped, abelian, and flows to the trivial fixed point
(\cref{comp:sb-verdict}).

\section{Objects and notation}
\label{comp:sb-notation}

The supplement adds its own symbol table for Parts B-I and B-II (\SB{app:base-backmatter}) and a local
table for the lattice part (\SB{tab:rg-lattice-notation}). The decoder covers what a reader meets first.

\begin{decoder}
Support \(\mathcal C_i\); design \(D=\{c_d\}\) & Contexts where agent \(i\) exists; a finite set of
sampled contexts, not a discretization of the base. & Local field support; probe points. \\
Relative frame \(h_i\) & Group element comparing belief and model frames, \(u_i^m=u_i^b\cdot h_i\). &
Bifundamental field of two gauge sectors. \\
Connections \(\omega_b,\omega_m\); transports \(\Omega_\gamma\) & Chosen rules for comparing fibers
along base curves. & Gauge potential and Wilson line. \\
Covariant defect \(\mathcal D\Phi\) & Failure of a belief-to-model cross map to commute with transport. &
\(\nabla\Phi\) for a Hom-valued field. \\
Graph flatness & Every closed walk of the agent graph has trivial link product. & Pure-gauge lattice
configuration. \\
Covariant jet \(D^\omega s\); pullback \(h^\omega_s\) & Rate of change of a belief field along the base
with frame rotation removed; its squared Fisher length. & \(D_\mu\phi\); gauged sigma-model metric. \\
Horizontal defect \(A_\Psi\) & Mismatch of fine and coarse connections under a coarse map over a base
map. & Mismatched gauge field under blocking. \\
Comparison data \(J_\ell\) & Declared identification of each blocked theory with a reference. &
Rescaling step of a Wilsonian RG map. \\
Relevant, marginal, irrelevant & Sign of a growth rate per unit log block scale, relative to declared
norms and normalizations. & Sign of an RG exponent \(y\). \\
\(z_e=\phi_x\psi_e\phi_{x+e_\iota}\) & Gauge-invariant link sign of the lattice model. & Unitary-gauge
link of the \(\mathbb Z_2\) Higgs model. \\
\(b=\tanh h\), \(t=\tanh j\) & Strengths of matter hopping \(h\) and plaquette coupling \(j\). &
High-temperature expansion variables. \\
Specification at every exterior & Coarse conditional laws defined and continuous at every boundary
condition, including null ones. & Gibbsianness of the renormalized measure. \\
Label \(\gamma\), bulk score \(\mathsf s_\gamma\) & A member of a finite catalog of lattice laws and its
extensive log-mass density. & Coexisting phases and their free energies. \\
\end{decoder}

Glyphs change meaning between parts. The letter \(b\) is the belief channel in Part B-I, a block ratio
\(b>1\) in Part B-II, and \(\tanh h\) on the lattice; \(t\) is a scale coordinate in
\SB{ch:wilsonian-limits} and \(\tanh j\) on the lattice; beta functions are \(\mathfrak b\), not the
attention rows \(\beta^x_{ij}\); on the lattice \(\rho=\tanh(h+2j)\) and \(\tau\) is a decay rate.
Words are narrower than a physicist expects. ``Pullback'' means the connection-relative pullback of
fiber tensors to the base, not Supplement G's restriction along a parameter map. ``Flat'' has four
meanings (\cref{comp:sb-flatness}). ``Connection'' can be a principal connection, the scale connection
of \SB{def:rg-scale-connection}, or a projected connection. ``Relevant'' is never a property of a flow
alone (\SB{def:rg-mode-exponents}), and ``fixed point'' is typed by tier: a fixed object at one tier need
not be fixed at another; six specific non-implications are proved in
\SB{prop:rg-fixed-object-nonimplication}.
``Declared'' means chosen rather than derived, and usually marks where modeling enters.

\section{Part B-I: base transport and multi-point designs}
\label{comp:sb-base}

\subsection{Frames, connections, and cross maps}
\label{comp:sb-frames}

Over a base, the belief and model channels of agent \(i\) may carry separately chosen local sections of
the one principal bundle, and the torsor structure makes their comparison canonical:
\[
u_i^m=u_i^b\cdot h_i,
\qquad
h_i\longmapsto g_i^bh_i(g_i^m)^{-1},
\qquad
\Theta_e^m=h_i^{-1}\Theta_e^bh_j\quad(e:j\to i).
\]
The relative frame \(h_i\) exists for any pair of representations, of any dimensions, and maps no
representation space into another (\SB{prop:geo-relative-frame}). It transforms on the left under
belief rechoices and on the right under model rechoices, like a bifundamental field. The two channels
carry one link on \(e\) exactly when the third equation holds (\SB{prop:geo-identified-links}); the main
manuscript's literal equality \(\Theta^b_e=\Theta^m_e\) is its gauge \(h_i=e_G\). The second atlas does
not enlarge the gauge group, whose active image at a context is still one copy of \(G\)
(\SB{prop:geo-two-atlas-active-image}); a genuine \(G_b\times G_m\) is an optional extension
(\SB{def:geo-product-extension}).

Connections are chosen data, with the textbook frame law \(A'=\operatorname{Ad}_gA-(\mathrm dg)g^{-1}\)
(\SB{prop:geo-connection-frame-laws}). The new object is the covariant defect of a cross map,
\[
(\mathcal D\Phi)_e(X)=T_e\Phi\bigl(H^b_eX\bigr)-H^m_{\Phi(e)}X,
\qquad\text{in the linear case}\qquad
D\Phi=\nabla^m\circ\Phi-\Phi\circ\nabla^b .
\]
The first term carries the belief connection's horizontal direction through \(\Phi\), the second is the
model connection's horizontal direction at the image. Their difference is vertical and vanishes exactly
when \(\Phi\) commutes with parallel transport along every base curve
(\SB{prop:geo-covariant-defect-parallelness}). At \(c_*\) there is no base vector and no defect. Across
scales, a coarse-graining map between principal bundles over a base map exists only when a
topological condition holds, as for the Hopf bundle, which no homomorphism \(z\mapsto z^n\) with
\(n\neq0\) maps into a trivial bundle (\SB{prop:rg-principal-scale-map-existence}).

\subsection{Four kinds of flatness}
\label{comp:sb-flatness}

At \(c_*\) every bundle is trivial and every connection vanishes, so ``flat'' can only mean graph
flatness. Over a base four conditions separate completely (\SB{prop:geo-four-notions}), with the first
decided by one Čech class common to both atlases (\SB{prop:geo-cech-class}).

\begin{center}
\small\renewcommand{\arraystretch}{1.15}
\begin{tabular}{@{}>{\raggedright\arraybackslash}p{0.2\textwidth}>{\raggedright\arraybackslash}p{0.3\textwidth}>{\raggedright\arraybackslash}p{0.42\textwidth}@{}}
\toprule
\textbf{Condition} & \textbf{Gauge-theory meaning} & \textbf{Separating witness} \\
\midrule
(a) Bundle triviality & A global gauge exists. & A flat \(S_3\)-bundle over a circle, glued by a
3-cycle, is nontrivial. \\
(b) Graph flatness & Agent-graph links are pure gauge. & Declared independently of (a), (c), (d). \\
(c) Connection flatness & \(F=0\). & \(x\mathrm dy\) on \(\R^2\): uniform field, not flat. \\
(d) Trivial holonomy on one loop & One Wilson loop is trivial. & \(x\mathrm dy\) around a circle
enclosing flux \(2\pi\); conversely an Aharonov--Bohm vortex is flat with nontrivial phase. \\
\bottomrule
\end{tabular}
\end{center}

For the program this means that flat links at \(c_*\), condition (b), say nothing about (a), (c), or
(d). The supplement also records that one principal class does not make the belief and model bundles
isomorphic (a cylinder and a Möbius band share one \(\{\pm1\}\)-bundle), and that a global section of an
associated bundle need not trivialize the principal bundle (\SB{prop:geo-cech-class}).

\subsection{When is a graph link a base transport?}
\label{comp:sb-graph-base}

When it is declared to be one. The standing hypothesis \SB{hyp:geo-graph-base-transport} chooses a
context per agent and a base curve per edge copy and requires
\(\widehat\rho_b(\Theta_e^b)=\Omega_{\gamma_e}\) and \(\widehat\rho_m(\Theta_e^m)=\widetilde\Omega_{\gamma_e}\).
Under it, graph loop products are base holonomies (\SB{prop:geo-graph-base-loop-transport}), a one-line
consequence of the homomorphism property. The supplement is explicit that the hypothesis is satisfiable
by definition, that no result realizes it for independently specified links
(\SB{open:rg-cover-graph-base-identification}), that recovering a connection from records is open
(\SB{open:rg-cover-holonomy-recovery}), and that even an exact realization concerns connection holonomy,
not bundle topology.

\begin{analogy}{structural}
At \(c_*\) the links are lattice gauge variables over a zero-dimensional base (\cref{comp:ch02}). The
hypothesis is Wilson's discretization read backwards: a lattice link is declared to be the path-ordered
exponential of a continuum potential along an assigned curve. In lattice gauge theory that is the
definition of the discretization, and here too it is a definition, not a theorem.
\end{analogy}

\subsection{Finite designs and shared links}
\label{comp:sb-designs}

A finite design is the multi-point version of \(\{c_*\}\) (\SB{def:prob-finite-design}). Under the
design-point product hypothesis, distinct design points are conditionally independent given the
structural data (\SB{hyp:gen-design-product}), and the free energy decomposes over the design
(\SB{cor:gen-design-factorwise-decomposition}); the hypothesis is substantive, since it excludes
correlation of an agent's states at two contexts beyond the structural data. If a submodel ties the
copies of one link across design points, each tied link ranging over all of \(G\), the surviving gauge
freedom is
\[
g_{d,i}^x=\bar g_i^xz_d,\qquad g_{d,j}^x=\bar g_j^xz_d,\qquad z_d\in Z(G),
\]
a context-independent rechoice times a central twist shared by both endpoints
(\SB{prop:gen-shared-link-stabilizer}). Tying a link across contexts turns local freedom along the design
into global freedom up to the center, which is \(\R^\times I\) for \(\GL(K,\R)\), \(\mathbb Z_N\) for
\(\mathrm{SU}(N)\), and trivial for \(\mathrm{SO}(3)\).

The finite construction supplies no probability law on a space of sections
(\SB{open:prob-continuum-theory}), and matching sections to finite laws at the design points determines
neither the joint nor the sections off the design (\SB{prop:prob-compatibility-nonidentifiability}). The
named obstruction is the one familiar from constructive field theory: reference measures are finite
products, and infinitely many non-probability factors give no measure, so a continuum theory must declare
a measure on fields directly. Until then every evidence-bound statement of the main manuscript and of
Supplement B is a statement about a finite design.

\subsection{Base-direction information geometry}
\label{comp:sb-pullback}

Once contexts vary, a belief field \(s\) changes along the base, and the connection separates the part of
that change that is only frame rotation:
\[
D^\omega s=\operatorname{ver}^\omega\circ Ts,
\qquad
h^\omega_s(X,Y)=g^F_{s(c)}\bigl(D^\omega s(X),D^\omega s(Y)\bigr)
\]
(\SB{def:pb-informational-pullbacks}). The first is the analogue of \(D_\mu\phi\); the second is the squared
Fisher speed of the belief along the base direction \(X\), whose square root is the Fisher--Rao distance per
unit displacement between the belief at \(c\) and the belief transported back from a neighbor. It is invariant under passive frame
changes (\SB{thm:pb-pullback-gauge-invariance}) but not under a change of connection: for
\(\omega'=\omega+\mathfrak d\),
\[
h^{\omega'}_s(X,Y)=h^\omega_s(X,Y)
+g^F\bigl(\mathsf J^s_{\mathfrak d}X,D^\omega sY\bigr)
+g^F\bigl(D^\omega sX,\mathsf J^s_{\mathfrak d}Y\bigr)
+g^F\bigl(\mathsf J^s_{\mathfrak d}X,\mathsf J^s_{\mathfrak d}Y\bigr),
\]
with \(\mathsf J^s_{\mathfrak d}\) the vertical field generated by the connection difference
(\SB{prop:pb-pullback-connection-change}). A belief constant over the base has zero pullback under the zero
connection and \(\mathrm c_1^2\mathrm dx^2\) under \(A'=\mathrm c_1\mathrm dx\), because the new connection
says the fiber is being shifted underneath it. This answers what a ``metric on the base built from beliefs'' is: a
metric relative to a chosen connection, and the documents choose none
(\SB{open:pb-canonical-pullback-geometry}). It is also only a semimetric, with radical the directions of
covariant constancy; constant rank gives a metric on a quotient vector bundle
(\SB{thm:pb-pullback-rank-quotient}), and a quotient manifold can fail even then, since the null directions
can form a contact distribution (\SB{prop:pb-contact-null-counterexample}). Combining the belief and model
channels needs declared weights (\SB{hyp:pb-weighted-product-geometry}).

Under a coarse map \(\Psi\) over a base map \(f\), fine and coarse base geometries obey the exact identity
\[
h^\omega_s-f^*\bar h^{\bar\omega}_{\bar s}=\delta_\Psi-\mathcal X_\Psi-\mathcal Q_\Psi
\]
(\SB{thm:pb-signed-base-comparison}). Here \(\delta_\Psi\) is the Fisher information lost through the fiber
channel, nonnegative for a Markov channel, the main manuscript's data-processing loss read along base
directions; \(\mathcal X_\Psi\) is a sign-indefinite cross term and \(\mathcal Q_\Psi\geq0\) a quadratic
term, both produced by the horizontal defect. With zero defect and a Markov channel the comparison is a
contraction and a
meta-agent's perceived base geometry is smaller than its children's
(\SB{cor:pb-base-fisher-contraction}, \SB{cor:pb-meta-perceived-geometry}); with nonzero defect it can be
negative at zero information loss (\SB{thm:pb-signed-positivity-criterion}). For maps induced by a principal
scale map the defect is the fundamental field of
\(\mathfrak A_{\mathcal P}=\mathcal P^*\bar\omega-\mathrm d\kappa\circ\omega\) and vanishes exactly when
\(\mathfrak A_{\mathcal P}\) takes values in the isotropy algebra of the coarse section value, with
\(\mathcal Q_\Psi\) the orbit Fisher form evaluated on \(\mathfrak A_{\mathcal P}\)
(\SB{thm:pb-isotropy-criterion}). Composed defects obey an ordered cocycle law
(\SB{thm:pb-anomaly-composition}), and the base Fisher defect composes up to an exact residual
(\SB{thm:pb-base-defect-cocycle}) whose relaxed pointwise problems are completely parameterized.

\begin{analogy}{structural}
The pullback is the kinetic metric \(\lvert D_\mu s\rvert^2\) of a gauged sigma model with the statistical
manifold as target, and its connection dependence is that of \(\lvert D_\mu\phi\rvert^2\) on \(A_\mu\). In
field theory the gauge field is dynamical or fixed by an equation of motion; the supplement supplies
neither. The isotropy criterion has the form of a Higgs mass term: a quadratic form in the gauge field
whose radical is the algebra fixing the vacuum value. The correspondence is of form only.
\end{analogy}

\subsection{Section-space histories}
\label{comp:sb-histories}

When an agent's state is a whole field, its inference history is a curve of sections, and it changes
pointwise vertically: at each context the belief moves within its fiber, and a base curve is not an
agent history (\SB{prop:hist-section-history-verticality}). On a strong Riemannian Hilbert manifold with
a free, proper, isometric gauge action, the quotient speed
\[
\lVert d\pi_Qv\rVert^2=\min_{\xi\in\mathfrak g}\mathsf g_Q(v-A_Q\xi,v-A_Q\xi)
\]
has a unique minimizer, so gauge-invariant durations are well defined
(\SB{thm:hist-strong-hilbert-quotient}); the minimization removes the velocity along gauge orbits, as a
Coulomb-type condition does. It fails for weak metrics: the \(L^2\) Fisher metric on \(H^1\) fields on a
circle gives a smooth proper quotient on which every tangent has zero speed
(\SB{prop:hist-proper-weak-fisher-quotient}), and that weak metric is itself the Fisher metric of a
natural observation family, so the obstruction is not an artifact of an exotic metric. Averaging a field
over the fibers of a base map loses a nonnegative amount of integrated Fisher information, given a Markov
fiber map, nonincreasing weights, and a jointly convex fiber chart; when the coarse metric is moreover
constant in the chart, the loss is exactly the sum of a channel loss, a weight gap, and a context variance
(\SB{thm:hist-averaging-defect}). Without joint convexity the conclusion fails: in the covariance sector of
the Gaussian chart the loss is negative on an explicit datum (\SB{prop:hist-averaging-convexity}).

\subsection{Covers, availability, and descent}
\label{comp:sb-covers}

\SB{ch:neighborhood-parents} is the part closest to the original picture of agents sprinkled over a base
with overlapping supports. Four agents on a square of contexts have strip supports that overlap in pairs
and triples with no common point (\SB{thm:rg-cover-nerve}), and each is present in a record with a smooth
availability probability \(w_i(c)\) that vanishes off its support (\SB{thm:rg-cover-availability}).
Attention then behaves as one would want over overlaps, with one qualification:
\(\beta^x_{ij}(c)=\eta^x_{ij}(c)/\alpha^x_i(c)\) is smooth where the receiver is present, while the event
law \(\eta\) is smooth everywhere and imposes no row where the receiver is absent
(\SB{prop:rg-cover-zero-occupancy}). The main manuscript's parent needs all its children present, so a
four-child parent here has empty domain; an availability parent weighting locally present children by
\(w_i/\sum w\) exists on the union of supports (\SB{thm:rg-cover-compatible-laws}). The weights are chosen,
and every compatibility statement holds by construction from one assumed ambient law
(\SB{hyp:rg-cover-ambient-law}); gluing independently supplied local laws is
\SB{open:rg-cover-local-gluing}.

The substantive result is the descent criterion: along a surjective submersion, a section descends
exactly when its image is constant on fibers, and a connection form exactly when its two pullbacks to the
fiber product agree (\SB{thm:rg-cover-descent-obstructions}). Around it, a coarse map with
context-dependent weights has defect \(O_{\mathrm c}(\mathrm dL)O_{\mathrm f}^{-1}\) even when all
connections agree (\SB{thm:rg-cover-nonflat-defect}); projecting a flat connection onto a moving quotient
creates curvature \(\bar{\mathsf g}^{-1}\mathrm{II}^\top\wedge\mathrm{II}\) from the second fundamental form,
and staged projections agree with the direct one when the intermediate quotient metric is propagated,
while resetting it can change the connection class (\SB{thm:rg-cover-projected-connection},
\SB{thm:rg-cover-quotient-composition}). Finally, a declared record
\(o_1=(\operatorname{Hol}_\iota-I)P_{\mathrm{pr}}Y+\varepsilon\) distinguishes two base-holonomy classes by a
frame-invariant statistic (\SB{thm:rg-cover-operational-holonomy}), the base companion of the main
manuscript's network-loop witness (\MT{thm:cg-loop-operational-holonomy}). The record reads the holonomy
by construction; the theorem shows that such a normalized experiment exists, not that the model's own
records are holonomy sensitive.

\begin{analogy}{exact}
The projected-connection curvature is the Gauss equation for an induced connection; with a flat ambient
connection and a moving rank-one image it is Berry curvature. The composition result says iterated
adiabatic projections agree with the direct one when each stage uses the metric induced by the last.
\end{analogy}

\subsection{How Part B-I connects to the base-manifold picture}
\label{comp:sb-vision}

Part B-I supplies the typed skeleton of the original picture: supports and Čech data, relative frames,
chosen connections, covariant defects, finite designs, a pulled-back base geometry, gauge-invariant
durations, availability-weighted attention over overlaps, and descent along base maps. It supplies no
dynamics for any base datum. The connection is chosen, not determined by inference; base measures and
weights are declared; the ambient law on a cover is assumed; and without a law on fields the evidence
bound over a base exists only on finite designs. The single-point graph links and the base transports
are related only by declaration. Part B-I is a precise list of what must be supplied when base transport
is taken up, not yet a theory of agents over a base.

\section{Part B-II: infinite networks and Wilsonian limits}
\label{comp:sb-wilson}

\subsection{What Wilsonian words require, and why finite depth has none}
\label{comp:sb-finite-depth}

The main manuscript coarse-grains finite networks by exact channels and keeps an effective action, a
closure residual, and finite-step multipliers (\cref{comp:ch11}). A Wilsonian step is ``block, then
identify the blocked theory with the original'', and the identification is extra structure: declared
comparison isomorphisms \(J_\ell\), giving the reference cocycle
\(\widehat{\mathfrak R}_\ell=J_{\ell+1}^{-1}\circ C_{\ell,\ell+1}\circ J_\ell\) (\SB{def:rg-comparison-data}).
Exponents are growth rates per unit of cumulative log block scale \(\varsigma_{n\leftarrow\ell}=\sum_k\log
b_k\),
\[
\chi_{n\leftarrow\ell}(u)=\frac{\log\lVert D_{n\leftarrow\ell}u\rVert_n-\log\lVert u\rVert_\ell}{\varsigma_{n\leftarrow\ell}},
\]
with \(D_{n\leftarrow\ell}\) the main manuscript's multipliers (\SB{def:rg-mode-exponents}). Where defined,
their canonical content is the forward Lyapunov filtration (\SB{prop:rg-lyapunov-filtration}); they are
invariant under tempered comparisons (\SB{thm:rg-tempered-comparison}), while a superexponential
comparison can change the apparent exponent (\SB{prop:rg-superexponential-distortion}); and no Oseledets splitting is
claimed, since a deterministic blocking sequence meets none of the hypotheses of the multiplicative
ergodic theorem \citep{Oseledets1968}. Even uniform projective contraction can leave an attracting
two-cycle rather than a fixed ray (\SB{prop:rg-contraction-no-fixed-ray}).

\SB{prop:rg-finite-depth-no-asymptotics} makes precise what every physicist knows about blocking a finite
system. With blocks of at least two agents the depth is at most \(\lfloor\log_2\lvert V_0\rvert\rfloor\), so
the limit defining an exponent does not exist; levels of different sizes admit no comparison isomorphism;
and exponent signs depend on the block schedule unless it is two-sided bounded: a mode doubled at each step has exponent 1 under
\(b_k=2\) and \(2/(n+1)\to0\) under \(b_k=2^{k+1}\). This does not conflict with renormalization at
\(c_*\) as a thermodynamic limit; it says that asymptotic vocabulary needs one, an infinite network or a
family in which size and depth grow together (\cref{comp:sb-universality}), whereas the main manuscript's
interaction multigraph is finite by definition (\MT{def:geo-graph-links}).

The exact step contributes no growth of its own: along an exact interaction orbit
\(\chi_{n\leftarrow\ell}(u)\leq\lvert V_n\rvert\log3/\varsigma_{n\leftarrow\ell}\), so no direction is
relevant in the native norms when \(\lvert V_n\rvert/\varsigma\to0\), and the same holds in the reference norm
when the comparisons are also tempered (\SB{cor:rg-interaction-no-relevant-mode}). The condition
\(\lvert V_n\rvert/\varsigma\to0\) is substantive: the bound grows with the number of coarse agents, and it
is vacuous for blocking by vertex ratios on a finite network. Relevance must come from a declared extensive normalization.
In the scalar Gaussian block each Hermite function of degree \(k\) has coefficient \(b^{1-k/2}\), so degree
one is relevant, two marginal, and higher degrees irrelevant, and the growth of degree one comes entirely
from the extensive replication normalization, of norm \(\sqrt b\), since the conditional expectation
contracts (\SB{def:rg-hermite-relevance}). Correlated coordinates
change this spectrum, a general \(\GL(d)\) action leaves the setting, and no link or holonomy enters
(\SB{prop:rg-hermite-scope}).

\begin{analogy}{structural}
In Wilson's renormalization group the blocking step, a conditional expectation, contracts; eigenvalues
above one come from the rescaling of fields that accompanies it \citep{wilson1974renormalization}. Part
B-II proves this division of labor for the exact step of the main manuscript. The Hermite spectrum is the
Gaussian fixed point of the central limit theorem read as a renormalization map, whose binary case is in
\citet{JonaLasinio2001}.
\end{analogy}

\subsection{Beta functions and fixed points}
\label{comp:sb-beta}

A beta function is a difference quotient on one reference space. For the action,
\(\mathfrak b^H_b[H;\varrho]=(\mathfrak R^H_b[H;\varrho]-H)/\log b\), and a change of reference
\(\varrho'=e^{-\Delta}\varrho\), \(H'=H-\Delta\), which leaves the weighted measure unchanged, shifts it by
the inhomogeneous term \(-\mathfrak b^H_b[\Delta;\varrho]\) (\SB{prop:rg-action-beta-reference-change}). A
beta function is thus a joint statement about the action and the reference trajectory. The retained beta
of a truncated family differs from the exact one by the transported closure residual over \(\log b\)
(\SB{prop:rg-retained-beta-residual}); a moving comparison frame gives even the identity step a nonzero
beta, and what survives smooth conjugacy are projective eigenvalue ratios at a fixed ray
(\SB{prop:rg-scheme-dependence}); and integer endpoints do not determine a continuous beta
(\SB{prop:rg-continuous-beta-underdetermined}).

A fixed theory is a measure pair, reference and weighted measure, fixed by the rescaled kernel. If the
kernel has a unique invariant law, every fixed pair has constant action
(\SB{prop:rg-fixed-pair-unique-ergodicity}), so a fixed theory with nontrivial action needs several
invariant laws; this depends on the typing choice that the reference is pushed by the same kernel. The
essential-spectrum theorem of the chapter has a hypothesis no exact arrow satisfies
(\SB{thm:rg-unital-essential-spectrum}), and one exact kernel, the circle doubling map, has spectral radius
1 on bounded functions and \(2^\alpha\) on Hölder functions (\SB{prop:rg-circle-norm-witness}), so a
spectrum belongs to a declared topology.

\subsection{The homogeneous Gaussian hierarchy}
\label{comp:sb-gauss-hierarchy}

On a homogeneous complete graph with equal blocks of size \(b\), common self term \(A\), and common
coupling \(W\), hard aggregation acts by
\[
A\longmapsto bA,\qquad W\longmapsto b^2W,
\]
informally, because a block has \(b\) self terms and two blocks have \(b^2\) cross couplings. On
\(\Sym^K\oplus\Sym^K\) the spectrum is \(b\) and \(b^2\), each with multiplicity \(K(K+1)/2\), so the ratio
outside the dominant sector is \(1/b\) (\SB{prop:rg-gaussian-dense-scaling}). Normalizing by \(b^2\) fixes the
coupling and contracts the self term toward a limiting ray with \(A=0\), a boundary operator whose
consensus zero modes make the Gaussian non-normalizable, so pinning, quotienting, or a small mass are
different infrared repairs. On general graphs the exponents follow the scaling of the cut: \(\xi=2\) for
the complete graph, \(\xi=0\) for a chain. This is mean-field power counting for Gaussian couplings.

It does not show attraction. Birkhoff's cone contraction cannot reach the boundary ray, since sum-only
aggregation is not primitive on the full cone and does not mix matrix directions
(\SB{prop:grg-full-cone-blocked}). Attraction is proved under one fixed primitive identification
\(B\otimes\operatorname{id}\) on edge indices:
\(r_B^{-\ell}(B^\ell W^{(0)})_f\to v_fM_\infty\) with \(M_\infty=\sum_eu_eW^{(0)}_e\) and \(u,v\) the Perron
vectors (\SB{conj:grg-fixed-b-attraction}). The spatial profile is attracted; the \(K\times K\) fiber matrix is
a Perron-weighted average of the initial couplings and is not selected. Whether any graph family admits such
an identification is \SB{open:rg-scalarized-attraction}, and nonautonomous identifications are not covered.
Attraction on an infinite hierarchy is therefore not established.

\subsection{Two-index limits, and what the absence of universality means}
\label{comp:sb-universality}

A thermodynamic scale statement has two indices, size \(n\) and depth \(\ell\), and three limits: fixed
depth after size, maximal depth, and a diagonal \(\ell(n)\) (\SB{def:rg-two-index-family}). They need not
agree: for the blocked path graph the fixed-depth limit of the spectral count is the arccos law, the
maximal-depth limit a single site (\SB{prop:rg-noncommuting-limits}). A countable index set carries a joint
law by Kolmogorov extension (\SB{thm:cg-projective-cylinder-law}), but consistency gives no ELBO limit and no
scale flow. That no universality class is identified (\SB{open:rg-two-index-limits}) means that the
supplement states only what one would need: a triangular family of finite agent networks at \(c_*\) with
consistency at fixed depth, uniform control of the coarse error, a diagonal scaling, and proved convergence
of a scheme-invariant observable (\SB{open:rg-two-index-limits}). Raw couplings and beta components are
scheme-dependent coordinates; a spectral exponent is a label only after convergence of the integrated
density of states is proved (\SB{def:rg-gaussian-spectral-labels}), and a finite graph has an atomic
spectrum. For the view of renormalization at \(c_*\) as a thermodynamic limit, the two-index family is the
typed home, and it is open.

\subsection{The numerical statements}
\label{comp:sb-numerical}

Supplement B has three NUMERICAL statements, all in Part B-II: a declared rescaling family on nested
six-agent cycles fails the semigroup law, with total variation 0.042 between direct and staged blocking
(\SB{rem:rg-compatibility-computation}); a frozen-attention projection scheme on a six-cycle and a
nine-cycle reaches fixed structures with vanishing pairwise block
(\SB{rem:rg-passive-blocking-computation}); and a row-regeneration scheme on the same instances keeps a
nonzero pairwise block (\SB{rem:rg-regenerated-attention-computation}). The schemes are declared
approximations, not the exact step, and the supplement says nothing here shows the exact pushforward
maintains interaction. The implementation has been retired and the values were not recomputed
(\SB{sec:rg-thermo-numerical-provenance}). The other rows of that index are exact rational checks of lattice
statements that rest on their proofs.

\section{Part B-III: the lattice cluster}
\label{comp:sb-lattice}

\subsection{The model and the flow}
\label{comp:sb-lattice-flow}

The source law lives on the edges of the square lattice, with matter signs \(\phi_x\), raw links \(\psi_e\),
and dressed signs \(z_e=\phi_x\psi_e\phi_{x+e_\iota}\). On a finite region with exterior \(\eta\) its weight
relative to counting measure is
\[
r^\eta_\Lambda(z)=\prod_{e\in\Lambda}(1+bz_e)\prod_{p\cap\Lambda\neq\varnothing}\bigl(1+t\chi_p(z\eta)\bigr),
\qquad b=\tanh h,\quad t=\tanh j
\]
(\SB{def:rg-spatial-source}). Since \(e^{hz}=\cosh h(1+bz)\) and \(e^{j\chi}=\cosh j(1+t\chi)\), this is
\(\exp(h\sum_ez_e+j\sum_p\chi_p)\), the \(\mathbb Z_2\) gauge theory of \citet{Wegner1971} coupled to a
\(\mathbb Z_2\) Higgs field as in \citet{FradkinShenker1979}, in unitary gauge: the first product is matter
hopping, a field on each link, and \(\chi_p\) is the plaquette holonomy. Small \(t\) is strong gauge
coupling, where \citet{OsterwalderSeiler1978} proved existence and analyticity. For \(0<b<1\) and \(0<t<1/6\)
the Gibbs law is unique and translation invariant, the matter signs are independent fair coins, and the
law depends on them only through \(z\) (\SB{thm:rg-spatial-source-law}).

The channel \(C_B\) keeps every \(B\)-th matter sign and replaces each coarse link by the product of the
\(B\) fine links along the straight path it covers; every fine link off the coarse grid lines is integrated
out. It is gauge equivariant, composes exactly, \(C_aC_B=C_{aB}\), and a coarse plaquette holonomy is the
product of the fine ones inside it (\SB{prop:rg-spatial-path-square}). The central theorem
(\SB{thm:rg-spatial-constrained-specification}) says that for \(0<b<1\), \(0<t<1/12\), and every \(B\), the
pushed law has a unique positive specification, defined and continuous at every coarse exterior including
null ones, whose dependence on the exterior beyond radius \(R\) decays like \(B^2(R+1)e^{-\tau(R+1)}\), and
that the domain \(12t<1\) does not depend on \(B\). The renormalized measure is Gibbsian, with a
renormalized Hamiltonian, at every block size on one fixed window, so the non-Gibbsianness catalogued by
\citet{VanEnterFernandezSokal1993} does not occur there. The generated interaction is complete with norm tending
to zero in \(B\) (\SB{thm:rg-spatial-generated-interaction}); for the integer factors \(a\in\{2,3,6\}\), on a
ball of character coefficients, the renormalization map \(\mathscr R_a\) is Fréchet \(C^1\) with contracting
linearization, so zero is its only fixed point there (\SB{thm:rg-spatial-local-operator}); and for a
compact box of rational source parameters and rational norm parameters, the source law enters that ball once
\(B\) exceeds an explicitly computable \(B_0\) (\SB{thm:rg-spatial-basin-entrance}).

The flow is to the trivial fixed point, and this should be read plainly. As \(B\) grows the blocked laws
converge on fixed cylinders to the fair product law, for every Gibbs law at any finite couplings
(\SB{prop:rg-spatial-cylinder-attraction}); informally, because a long Wilson-line parity multiplies many
signs whose conditional means stay away from \(\pm1\). Interacting fixed points, critical exponents, and a continuum
gauge-field measure lie outside the construction (\SB{open:rg-spatial-generality}). What is nontrivial is the
control: Gibbsianness uniform in block size, a complete interaction with explicit bounds, and a
differentiable map, all constructed. At each fixed depth the Fisher information per coarse cell is a
\(C^2\) density that is a regular metric on a nonempty open region (\SB{thm:rg-spatial-fisher-density}), and
\(\mathcal I_{aB}\preceq a^2\mathcal I_B\): per fine cell information decreases, the data-processing
inequality (\SB{prop:rg-spatial-fisher-contraction}).

\begin{analogy}{exact}
The source is the \(\mathbb Z_2\) gauge--Higgs model in unitary gauge, and the straight-path channel is
blocking of gauge links by straight Wilson lines along the coarse grid, a product map rather than
decimation.
\end{analogy}

\begin{analogy}{structural}
The open-ball construction is of the kind \citet{Israel1981} gave for decimation and Kadanoff maps of
Ising-spin interactions at high temperature; here the channel and the interaction space are different, and
the supplement separates its results from those maps.
\end{analogy}

\subsection{Pressure, label selection, and the specific VFE}
\label{comp:sb-pressure}

\SB{ch:geometry-pressure} places a finite catalog of such laws on one lattice; each label is a value of the
structural mark of the main manuscript's partition catalog (\MT{def:rg-geometry-ensemble},
\cref{comp:ch12}). Pressures exist with explicit surface rates, \(\lvert\mathfrak p-n^{-2}\log
Z_n\rvert\leq4\mathsf m_1/n\) (\SB{thm:rg-pressure-boundary}); existence is standard \citep{Georgii1988} and
the rates are the addition. Under global normalization the label with the largest bulk score takes all the
weight exponentially fast in the volume, and at a tie the outcome is decided by subextensive terms, not by
pressures (\SB{thm:rg-pressure-selection}). Forgetting the label can destroy Gibbsianness outright: a mixture
of two Bernoulli product laws has no continuous conditional at any exterior (\SB{thm:rg-geometry-erasure}).

Observation is handled for one record: every edge sensor reports \(+\)
(\SB{hyp:rg-pressure-positive-record}). It shifts each branch field from \(b_\gamma\) to
\((b_\gamma+\eta_\gamma)/(1+b_\gamma\eta_\gamma)\), and the specific VFE of a label-retaining recognition
satisfies \(\min_{\mathbb Q}\mathfrak f=-\bar e\), attained on the maximal posterior-score labels with the
shifted Gibbs laws (\SB{prop:rg-pressure-vfe-minimum}): the Gibbs variational principle in the
thermodynamic limit, with a finite label. Along a fixed blocking word
\(B^2\mathfrak f_{\mathrm{fine}}=\mathfrak f_{\mathrm{coarse,induced}}+\mathfrak d_{\mathrm{fib}}\) with
\(\mathfrak d_{\mathrm{fib}}\geq0\) (\SB{thm:rg-pressure-induced-vfe}), the volume-limit form of the main
manuscript's exact coarse free-energy identity. A general record produces random fields of both signs;
its quenched evidence density is \SB{open:rg-pressure-quenched-record}.

\subsection{Thermodynamic Fisher geometry}
\label{comp:sb-thermo-fisher}

On compact source boxes with \(h>0\) and \(0<j<\operatorname{atanh}(1/12)\), the native pressure is \(C^4\),
and since the native law is an exponential family in \(\mathsf g=(h,j)\) its Fisher density is the Hessian
of the log-partition density, \(\tfrac1NI\to\mathsf H=\nabla^2_{\mathsf g}\mathsf A\) in \(C^2\)
(\SB{thm:rg-thermo-native-cfour-pressure}, \SB{prop:rg-thermo-native-ctwo-metric}); pushes along a fixed word
have the same \(C^2\) density per coarse cell as the infinite-law marginals
(\SB{thm:rg-thermo-native-free-push-ctwo}). The main manuscript's finite-mixture asymptotics
(\MT{sec:cg-retained-label-asymptotics}) then hold with \(N=L^2\) cells
(\SB{cor:rg-thermo-native-label-limits}): with a strict winner the Fisher density converges to the winner's;
at a tie \(I/N^2\to\operatorname{Cov}_w(\mathsf v_\gamma)\), the covariance of the bulk-score slopes, which is
information about the label; and in a window of width \(1/N\) the label law is logistic.

\begin{analogy}{exact}
For an exponential family the Fisher metric is the Hessian of the log-partition function, the thermodynamic
metric of fluctuation theory \citep{Crooks2007}, per cell the susceptibility matrix.
\end{analogy}

\begin{analogy}{structural}
The tie and window statements have the form of finite-size rounding at a first-order transition, which
\citet{BorgsKotecky1990} derive from the lattice model itself. Here the phase mixture is prescribed through the
label weights, which makes the statement elementary; the correspondence is one of form.
\end{analogy}

\subsection{The typing question and the thermodynamic-limit reading}
\label{comp:sb-typing}

No lattice statement declares what a site is (\SB{open:rg-lattice-typing}). In the first reading the sites
are agents at \(c_*\), \(G=\mathbb Z_2\), the raw link is an independent link, the dressed sign is the per-edge
twist, and a plaquette is the holonomy of an agent four-cycle. In the second the sites are points of a
discretized base, the raw link is the transport of a \(\mathbb Z_2\) connection, a plaquette is a discrete
curvature, and the channel is a scale map over a base map. Every statement holds under either reading.

The first reading is the reader's own: links as lattice gauge theory over a zero-dimensional base, and
renormalization at \(c_*\) as a thermodynamic limit. Under it the lattice cluster is the one place in the
three documents where a thermodynamic limit of a link law on an agent graph is actually constructed. The
open problem
states what is missing before the reading can be adopted: the lattice law must be the link marginal of a
normalized joint with belief and model channels, a recognition law, and an ELBO on a countably infinite
agent set, whereas the main manuscript's multigraph is finite and the lattice family carries none of these.
The lattice part also fixes every blocking word before its volume limit; joint growth of depth and volume
is part of \SB{open:rg-spatial-generality}, and where depth and volume do grow together the limit is the
trivial one. Even under the first reading the group is finite and abelian, so the per-edge twist is gauge
invariant and nonabelian holonomy classes never appear, far from the target of 512 or more agents with
\(\GL(K,\R)\), \(\GL(K,\mathbb C)\), or \(\mathrm{SU}(N)\) links; the domain is a single-phase strong-coupling
region, with a unique Gibbs law and a smooth pressure; and observation enters through one constant record. The lattice cluster proves that the
coarse-channel calculus can be run to completion in infinite volume for one abelian family. It is not
evidence about formation in the target class.

\subsection{What is new, and the classical antecedents}
\label{comp:sb-antecedents}

The model and tools are classical: \citet{Wegner1971}, \citet{FradkinShenker1979}, and
\citet{OsterwalderSeiler1978} for the model; Dobrushin uniqueness \citep{Dobrushin1968} and boundary
comparison \citep{Georgii1988}; Kolmogorov extension; Kozlov's Gibbs representation \citep{Kozlov1974};
Israel's high-temperature renormalization maps \citep{Israel1981}; and \citet{VanEnterFernandezSokal1993}
for Gibbsian pathologies, of which \SB{thm:rg-geometry-erasure} is a Bernoulli case in sharper form. The
principal candidate for a new result is the block-size uniformity of \(12t<1\) in
\SB{thm:rg-spatial-constrained-specification}: the surveyed high-temperature results concern decimation,
Kadanoff, and averaging maps of Ising-spin interactions, not this gauge--matter path channel, and the proof
avoids a block estimate whose constant grows like \(Bt\). The literature check is limited to the works the
chapter cites, and priority is not claimed.

\section{What is proved, and how solid it is}
\label{comp:sb-scorecard}

Established no-go statements and separating counterexamples count as solid. The scope column says how far
each result's generality is from the program's target.

\begin{scorecard}
\SB{prop:geo-relative-frame}, \SB{prop:geo-identified-links} & A unique relative frame; one shared link
means \(\Theta^m_e=h_i^{-1}\Theta^b_eh_j\). & \SOLID & Standard in substance (torsor). \\
\SB{prop:geo-four-notions} & Bundle triviality, graph flatness, connection flatness, and one-loop holonomy
are independent. & \SOLID & \(\mathrm U(1)\) and \(S_3\) witnesses; at \(c_*\) (a), (c), (d) hold trivially and
only (b) has content. \\
\SB{hyp:geo-graph-base-transport} & A graph link is a base transport only by declaration. & \CHOICE &
Realization for independent links is open. \\
\SB{prop:gen-shared-link-stabilizer} & Tying a link across contexts leaves constant rechoices times a
central twist. & \SOLID & Admissible rechoices form a groupoid. \\
\SB{open:prob-continuum-theory} & No law on fields; evidence bounds are finite-design statements. & \OPEN &
No infinite product of reference measures. \\
\SB{thm:pb-pullback-gauge-invariance}, \SB{prop:pb-pullback-connection-change} & The base Fisher pullback is
passively covariant and changes with the connection by exact terms. & \SOLID & No connection-free base
metric; only a semimetric. \\
\SB{thm:pb-signed-base-comparison} & Fine minus coarse base Fisher is channel loss minus two anomaly
terms. & \SOLID & Exact; can be negative at zero loss. \\
\SB{thm:pb-isotropy-criterion} & The horizontal defect vanishes exactly when the defect form lies in the
isotropy algebra. & \SOLID & Pointwise; no compatible connection claimed. \\
\SB{thm:rg-cover-compatible-laws} & An ambient law gives compatible local data; availability parents glue.
& \COND & Holds by construction (functoriality) from one assumed ambient law; gluing from local data
open. \\
\SB{thm:rg-cover-descent-obstructions} & Sections and connections descend exactly under fiber constancy
and pullback agreement. & \SOLID & Standard in substance. \\
\SB{thm:rg-cover-operational-holonomy} & A declared record distinguishes two base-holonomy classes. &
\TOY & Base-loop instance of the MT witness, whose mechanism is general (\SOLID{} in \cref{comp:ch10}); one
loop, one connection family, Gaussian record; recovery beyond it is \SB{open:rg-cover-holonomy-recovery}. \\
\SB{prop:rg-finite-depth-no-asymptotics} & A finite network has bounded depth, no comparisons, and
schedule-dependent exponents. & \SOLID & A no-go for finite \(N\), not for \(N\to\infty\). \\
\SB{cor:rg-interaction-no-relevant-mode} & The exact step has no relevant direction; relevance needs an
extensive normalization. & \COND & Needs \(\lvert V_n\rvert/\varsigma\to0\): the bound is exponential in
\(\lvert V_n\rvert\), so it is silent when coarse levels keep many agents, and vacuous for finite vertex-ratio
blocking. Tempered comparisons are needed only for the reference-norm clause. \\
\SB{prop:rg-gaussian-dense-scaling} & Complete graph: \(A\mapsto bA\), \(W\mapsto b^2W\); dense
normalization contracts the self term by \(1/b\). & \SOLID & Exact; limiting ray is a boundary operator. \\
\SB{conj:grg-fixed-b-attraction} & Under a fixed primitive identification orbits converge; the fiber matrix
is not selected. & \COND & Standard in substance (Perron--Frobenius on a primitive \(B\)); existence of
the identification is open. \\
\SB{rem:rg-compatibility-computation}, \SB{rem:rg-passive-blocking-computation},
\SB{rem:rg-regenerated-attention-computation} & A rescaling family fails the semigroup law; frozen attention
kills the pairwise block, regeneration keeps it. & \NUM & Six- and nine-agent schemes; not
recomputed. \\
\SB{open:rg-two-index-limits} & Joint size and depth limits and universality for agent networks. & \OPEN &
Limits need not commute. \\
\SB{thm:rg-spatial-constrained-specification} & Every straight-path blocking of the \(\mathbb Z_2\)
gauge--matter law at \(12t<1\) is Gibbsian at every exterior, on a domain independent of block size. & \SOLID & Principal candidate for a new result; \(\mathbb Z_2\), strong
coupling. \\
\SB{thm:rg-spatial-generated-interaction}, \SB{thm:rg-spatial-local-operator} & Complete interaction tending
to zero; for \(a\in\{2,3,6\}\), a \(C^1\) map on a ball whose only fixed point is zero. & \SOLID & Israel-type
in form, different channel; trivial fixed point. \\
\SB{prop:rg-pressure-vfe-minimum} & Specific VFE minimum is minus the evidence density. & \COND & Only for the
constant all-positive record. \\
\SB{cor:rg-thermo-native-label-limits} & \(C^2\) Fisher densities; strict, tie (\(N^2\)), and \(1/N\)-window
limits. & \SOLID & Compact boxes; blocking word fixed before volume. \\
\SB{open:rg-lattice-typing} & Are lattice sites agents at \(c_*\) or base points? & \OPEN & No lattice
statement is yet about agents. \\
\end{scorecard}

\section{Conditions that matter}
\label{comp:sb-hygiene}

\begin{hygiene}
Substantive conditions, which could fail in the cases the program cares about: chosen connections
(\SB{def:geo-connections}), relative to which every base-direction statement is made; the declaration
\SB{hyp:geo-graph-base-transport}; conditional independence of design points (\SB{hyp:gen-design-product});
declared base measures, channel weights, and availability weights; constant rank of the covariant jet
(which can jump, as in the rank-jump witness of \SB{prop:pb-contact-null-counterexample}), integrable null directions, and basic tensors for quotient
geometries; chart joint convexity (\SB{hyp:hist-joint-convexity}),
which fails for Gaussian covariances; a strong metric on a declared manifold of fields
(\SB{hyp:hist-regular-section-space}), which the Fisher metric of the weak-metric witness is not; one
ambient law on a cover
(\SB{hyp:rg-cover-ambient-law}); declared comparison data, infinite blocking sequences, and tempered
comparisons for every asymptotic statement; one fixed primitive identification for attraction; and on the
lattice the group \(\mathbb Z_2\), small plaquette coupling with positive field, compact source boxes,
blocking words fixed before volume limits, controlled label growth (\SB{hyp:rg-pressure-label-growth}), and
the constant all-positive record (\SB{hyp:rg-pressure-positive-record}).

Hygiene a physicist can take for granted on a first reading: standard Borel spaces, measurability of section
responses (\SB{hyp:prob-context-section-response}), smoothness classes \(C^1\) to \(C^4\), paracompactness and
the constant-rank bookkeeping that restricts the defect classifications to regular strata (a technical
restriction, unlike the rank of the jet itself), integrability of the integrals, the submersion hypothesis of descent,
boundary conventions, and centering of squares in local limits.
\end{hygiene}

Toy content is limited and labeled: the four-agent cover, the holonomy witness, the scalar Hermite spectrum,
and the numerical six- and nine-agent schemes. The many small counterexamples of Part B-I are not toys;
each refutes a general implication and is a solid negative result. The lattice is infinite but special in
its group, abelian and finite, and in its single-phase coupling region. Its mechanism, keeping a renormalization map
Gibbsian uniformly in block size by a transport argument, is specific to strong coupling and unshown for
nonabelian or non-compact links; its constants (\(12\), \(\mathrm q\), \(\mathrm c_1\)) say nothing about the
target class.

\section{Open ends}
\label{comp:sb-open}

The supplement lists twenty open problems (\SB{tab:base-open}). For Part B-I, the unknowns are a probability
law on fields with continuum limits of the bound (\SB{open:prob-continuum-theory}); the orbit space of the
two-atlas coordinate action (\SB{open:gen-gauge-quotient}); a base geometry fixed by the model rather than a
chosen connection (\SB{open:pb-canonical-pullback-geometry}); whether declared recognition families have the
structure the history theorems use (\SB{open:hist-section-tier-admission}) and whether descending fields form
a smooth submanifold (\SB{open:hist-projectable-set}); gluing local laws on a cover
(\SB{open:rg-cover-local-gluing}); realizing independent graph links as base transports
(\SB{open:rg-cover-graph-base-identification}); recovering a connection from records
(\SB{open:rg-cover-holonomy-recovery}); quotients across rank changes (\SB{open:rg-cover-quotient-extension});
and informative averaging along a base map (\SB{open:rg-cover-base-aggregation}). Each is settled by a
construction for a declared class or by an obstruction.

For Part B-II, the unknowns are rescaling kernels on agent networks obeying the semigroup law
(\SB{open:rg-rescaling-compatibility}); a multiplicative ergodic structure for scale cocycles
(\SB{open:rg-multiplicative-ergodic}); relevance spectra of correlated or gauge-mixed Gaussian blocks
(\SB{open:rg-hermite-extensive-relevance}) and a basin for the Gaussian block law
(\SB{open:rg-hermite-nonlinear-attraction}); a canonical continuous beta (\SB{open:rg-continuous-beta}); joint
size and depth limits with a convergent scheme-invariant observable (\SB{open:rg-two-index-limits}); and a
fixed primitive identification for some graph family (\SB{open:rg-scalarized-attraction}). For Part B-III,
they are the typing (\SB{open:rg-lattice-typing}), settled by proving one reading; critical points,
depth-uniform Fisher floors, joint depth and volume limits, and self-consistent curvature actions
(\SB{open:rg-spatial-generality}); and the quenched evidence density for random records
(\SB{open:rg-pressure-quenched-record}), which needs every-exterior control of random fields of both signs.

Three gates carry most of the weight. Without \SB{open:prob-continuum-theory} there is no evidence bound over a
base. Without \SB{open:rg-cover-graph-base-identification} the agent graph's links and the base's connection
stay separate objects joined only by declaration. And \SB{open:rg-two-index-limits}, with
\SB{open:rg-lattice-typing}, gates renormalization as a thermodynamic limit: the typed statement exists, the
one constructed instance is an untyped abelian lattice at strong coupling, and no agent-network instance with
its recognition law exists.

\section{Verdict}
\label{comp:sb-verdict}

\begin{verdictbox}{\NARROWER}
\textbf{What it says it does.} It collects what the theory requires once contexts vary, networks become
infinite, and agents sit on a lattice, stating for each extension which single-point statements survive and
which additional data it consumes (\SB{ch:base-introduction}). It declines to identify the base with physical
space, graph links with base transports, or any continuum law, and declines asymptotic vocabulary for finite
networks.

\textbf{What the statements give.} Part B-I gives the kinematics of a base extension, with exact
identities and separating counterexamples: relative frames, covariant defects, four independent flatness
notions, the shared-link stabilizer, an exact fine-to-coarse base Fisher comparison with a sharp isotropy
criterion, durations on strong metrics, and descent. Part B-II gives a typing of Wilsonian vocabulary, no-go
and scheme-dependence results, exact Gaussian power counting, and one conditional attraction theorem. Part
B-III gives a complete constructive renormalization of one \(\mathbb Z_2\) family at strong coupling, with
pressure, label selection, specific VFE, and Fisher densities in the volume limit.

\textbf{The gap.} ``Base transport'' is transport by a chosen connection, not transport derived from
inference; base measures, weights, and ambient laws are declared; no law on fields exists; and graph links
become base transports only by declaration. ``Infinite networks'' names the apparatus an infinite network
would need; no infinite agent network with channels, recognition law, and ELBO is constructed, no exponent or
fixed point of the agent flow is established, and universality is a stated program. The lattice cluster is
untyped, abelian, confined to strong coupling, observed through one record, and flows to the trivial fixed
point. Each part does what its opening paragraph says; the supplement as a whole is a precise map of what base
transport, infinite networks, and spatial realization would require, not a theory of any of them.
\end{verdictbox}
